%!TEX root=../main.tex
\chapter{Introduction}
\chaptermark{Introduction}
\label{ch:introduction}

\section{Motivation}
Multi-robot systems increasingly operate in settings where wireless
communication cannot be assumed to be instantaneous or reliable --
warehouses with dense metal shelving, underground or underwater
deployments, and outdoor sites with variable network coverage all
introduce meaningful communication delay. Most trajectory planners for
multi-robot coordination assume each robot has an up-to-date view of its
neighbours' planned paths \citep{VanDenBerg2011}; under delay, that
assumption breaks down and the planner's safety guarantees no longer hold.

\section{Problem Statement}
We consider a fleet of robots planning trajectories in a shared workspace,
where communication between any pair of robots is delayed by an unknown
but bounded amount. Each robot must plan a trajectory that remains
collision-free with respect to its neighbours' actual (not merely
last-known) positions, without assuming synchronized clocks or a
zero-delay communication channel \citep{Zhou2017}.

\section{Major Contributions}
Our main contributions are as follows:
\begin{itemize}
  \item A trajectory planning method that reserves a time-varying safety
    margin around each robot's path, sized according to a known
    communication delay bound rather than a fixed worst-case margin.
  \item A simulation study comparing the method against a delay-unaware
    baseline across a range of delay magnitudes and robot densities.
  \item A physical validation on a fleet of ground robots under
    artificially injected communication delay.
\end{itemize}

\section{Outline of the Thesis}
\cref{ch:introduction} introduces the problem and motivates the need for
delay-aware planning. \cref{ch:related} reviews prior work on multi-robot
trajectory planning and delay-tolerant coordination. Later chapters
(omitted from this starter document) would present the planning method,
simulation results, and physical robot trials in detail.
